\documentclass[a4paper,12pt]{report}
\setcounter{secnumdepth}{5}
\setcounter{tocdepth}{3}
\newcounter{ZhRenew}
\setcounter{ZhRenew}{1}
\newcounter{SectionLanguage}
\setcounter{SectionLanguage}{1}
\input{/usr/share/latex-toolkit/template.tex}
\begin{document}
\title{Combinatorics}
\author{沈威宇}
\date{\temtoday}
\titletocdoc
\section{Combinatorics (組合數學)}
\subsection{Counting principle (計數原理) or Fundamental principle of counting (基本計數原理)}
\sssc{Rule of sum or addition principle (加法原理)}
If there are $a$ possible outcomes for an event (or ways to do something) and $b$ possible outcomes for another event (or ways to do another thing), and the two events cannot both occur (or the two things can't both be done), then there are $a+b$ total possible outcomes for the events (or total possible ways to do one of the things).
\sssc{Rule of product or multiplication principle (乘法原理)}
If there are $a$ possible outcomes for an experiment (or ways to do something) and $b$ possible outcomes for another experiment (or ways to do another thing), and the two experiments are independent (or the two things must both be done), then there are $a\cdot b$ total possible outcomes for the sequential experiment of the two experiments (or total possible ways to do both things).
\sssc{Counting principle for power sets (冪集計數原理)}
\[\abs{2^A}=2^{\abs{A}}\]
\sssc{Principle of inclusion-exclusion or inclusion–exclusion principle, PIE (排容原理/取捨原理)}
For indexed family $(A_i)_{i=1}^n$ with $n\in\bbN$,
\[\abs{\bigcup_{i=1}^nA_i}=\sum_{J\subseteq\bbN_{\leq n}\land J\neq\varnothing}(-1)^{|J|+1}\abs{\bigcap_{j\in J}A_j}.\]
\begin{proof}
When $n=1$,
\[|A_1|=|A_1|.\]
When $n=2$,
\[\abs{A_1\cup A_2}=\abs{A_1\setminus(A_1\cap A_2)}+\abs{A_2\setminus(A_1\cap A_2)}+\abs{A_1\cap A_2}=\abs{A_1}+\abs{A_2}-\abs{A_1\cap A_2}.\}
Assume that is holds for $n\in\mathbb{N}_{\ge 2}$. For $n+1$,
\[\ba
\abs{\bigcup_{i=1}^{n+1}A_i}&=\abs{\bigcup_{i=1}^nA_i}+\abs{A_{n+1}}-\abs{\qty(\bigcup_{i=1}^nA_i)\cap A_{n+1}}\\
&=\sum_{i=1}^n\abs{A_i}-\sum_{1\le i<j\le n}\abs{A_i\cap A_j}+\sum_{1\le i<j<k\le n}\abs{A_i\cap A_j\cap A_k}-\cdots+(-1)^{n+1}\abs{\bigcap_{i=1}^nA_i}+\abs{A_{n+1}}-\qty(\sum_{i=1}^n\abs{A_i\cap A_{n+1}}-\sum_{1\le i<j\le n}\abs{A_i\cap A_j\cap A_{n+1}}+\sum_{1\le i<j<k\le n}\abs{A_i\cap A_j\cap A_{n+1}}-\cdots+(-1)^{n+1}\abs{\qty(\bigcap_{i=1}^nA_i)\cap A_{n+1}})\\
&=\sum_{i=1}^{n+1}\abs{A_i}-\sum_{1\le i<j\le n+1}\abs{A_i\cap A_j}+\sum_{1\le i<j<k\le n+1}\abs{A_i\cap A_j\cap A_k}-\cdots+(-1)^{n+2}\abs{\bigcap_{i=1}^{n+1}A_i}
\ea\]
By mathematical induction, the statement is true for all $n\in\mathbb{N}$.
\end{proof}
\sssc{Pigeonhole principle (鴿籠原理/鴿洞原理/鴿巢原理) or Dirichlet's drawer principle (狄利克雷抽屜原理)}
Given $n,k\in\bbN$, if $k$ pigeons fly into $n$ cages, then there must be at least one cage containing at least $\lceil\frac{k}{n}\rceil$ pigeons.
\sssc{Generalized pigeonhole principle}
Given $n,k_1,k_2,\ldots,k_n\in\bbN$, if $\sum_{i=1}^nk_i-n+1$ pigeons fly into $n$ cages, then there exists at least one $i\in\bbN_{\le n}$ such that the $i$th cage contains at least $k_i$ pigeons.
 \subsection{Factorial (階乘)}
\sssc{Factorial}
\[n!\coloneq\prod_{i=1}^ni,\quad n\in\bbN_0.\]
\sssc{Falling factorial}
\[(n)_m\coloneq\prod_{i=0}^{m-1}(n-i),\quad n\in\bbC\land m\in\bbN.\]
\ssc{Gamma function (伽瑪函數)}
\sssc{Definition}
\[\Gamma(z)\coloneq\int_0^{\infty}t^{z-1}e^{-t}\dd{t},\quad\Re(z)>-1.\]
The integral diverges if $\Re(z)\leq -1$.
\begin{proof}
PLACEHOLDER
\end{proof}
\sssc{Recurrence}
\[\Gamma(z+1)=z\Gamma(z),\quad\Re(z)>-1.\]
\begin{proof}
PLACEHOLDER
\end{proof}
For $z\in\bbN$,
\[\Gamma(z)=(z-1)!.\]
For $m\in\bbN\land z\notin\bbN_{\le m}$,
\[\frac{\Gamma(z)}{\Gamma(z-m)}=(z-1)_m.\]
\subsection{直線排列}
\sssc{盡相異物排列}
從$n$相異物中取$m$個排成一直列（$0\leq m\leq n$），正逆序視為二，其排列總數$P_m^n=\frac{n!}{\qty(n-m)!}$。
\sssc{不盡相異物排列}
設有$n$物，共有$k$種，第$i$種有$m_i$個。全取排成一直列，正逆序視為二，其排列總數為$\frac{n!}{\prod_{i=1}^n m_i!}$。
\sssc{盡相異物重複排列}
從$n$類相異物中，任取$m$個排成一直列，正逆序視為二，其中每類物品的個數均不小於$m$且可重複選取，其排列總數為$n^m$。
\sssc{錯排（Derangement）}
$n$相異物全取作直線排列，其中$m$物（$m\leq n$）依次被限制不能排列於相異單一指定位置之排列，其方法數稱$D_m^n$；當$m=n$，錯排方法數稱$D_n$。
\subsubsection{全錯排}
\tb{遞迴式}：
\[
\begin{cases}
D_0 &= 1 \\
D_1 &= 0 \\
D_n &= (n-1) \cdot (D_{n-1}+D_{n-2})\quad\text{（當\protect\ }n > 2\text{）}
\end{cases}
\]
\tb{一般式}：
\[
\begin{aligned}
D_n &= n! \cdot \left(1+\sum_{i=1}^n \frac{(-1)^i}{i!} \right) \\
&= \sum_{i=0}^n (C_i^n \cdot (n-i)! \cdot (-1)^i
\end{aligned}
\]
\tb{錯排機率}：
\[ \lim_{n \to \infty} \frac{D_n}{n!} = \frac{1}{e} \]
\subsubsection{非全錯排}
\[ D_m^n = \sum_{i=0}^m \left((-1)^i \cdot C_i^m \cdot (n-i)!\right) \]
\subsection{環狀排列（Cyclic permutation）}
環狀排列：$n$物全取排列成環狀，不同方法之判定僅考慮相對位置，不考慮絕對位置，惟翻轉（順時針$\rightleftharpoons$逆時針）視為二種。
\subsubsection{盡相異物環狀排列}
\[ \frac{n!}{m\cdot (n-m)!} \]
\subsubsection{不盡相異物環狀排列}
$n$ 不盡相異物，全取作環狀排列，求其方法數。
\begin{itemize}
\item \textbf{法一：}\\
子循環：一個長度$n$的給定直線排列，若其前$\frac{n}{m}$物重複$m$次等同於原排列，且不存在$>1$的$k$使原排列的前$\frac{n}{m\cdot k}$物重複$k$次等同於原排列的前$\frac{n}{m}$物，則稱原排列的前$\frac{n}{m}$物為一個長度$\frac{n}{m}$的子循環，稱原排列之子循環長度為$\frac{n}{m}$、子循環數目為$m$。 \\
令該$n$不盡相異物的所有可能直線排列中，有$d_i$個之子循環長度為$\ell_i$，且$\sum_{i=1}^m d_i=$該$n$不盡相異物直線排列方法數，則：
\[ \text{所求}=\sum_{i=1}^m \frac{d_i}{\ell_i} \]
\item \textbf{法二：}\\
令該$n$不盡相異物可分為$k$相異類，第$i$類（$1\leq i\leq k$）有$m_i$件相同物。最大公因數$\gcd(m_1, m_2, m_3, \ldots, m_k)=g$。令所求為$R$。
\begin{itemize}
\item 最大公因數$g$為$1$：
\[ R=\frac{\left(n-1\right)!}{\prod_{i=1}^k\left(m_i!\right)} \]
\item 最大公因數$g$為一質數$p$：
\[ R=\frac{1}{n}\cdot \left(\frac{n!}{\prod_{i=1}^k\left(m_i!\right)}-\frac{n!}{\prod_{i=1}^k\left(\frac{m_i}{p}!\right)}\right)+\frac{\left(\frac{n}{p}-1\right)!}{\prod_{i=1}^k\left(\frac{m_i}{p}!\right)} \]
\item 最大公因數$g$為任一正整數：
\[
\begin{aligned}
\text{令\protect\ }g & \text{之所有相異質因數由小到大依序為}p_1, p_2, p_3, \ldots, p_q。 \\
R &= \sum_{j=0}^q \left((-1)^j \cdot \sum_{\left|Y\right|=j} \frac{\left(\sum_{i=1}^k \frac{m_i}{\prod_{y \in Y} y}\right)!}{\prod_{i=1}^k \left(\frac{m_i}{\prod_{y \in Y} y}!\right)}\right) \\
& \text{，其中 } Y \subseteq \{p_1, p_2, p_3, \ldots, p_q\} \\
& \text{，定義 } \prod_{y \in \varnothing} y = 1
\end{aligned}
\]
\end{itemize}
\end{itemize}
\ssc{Choose with replacement}
Choose with replacement refers to a selection process where each item can be selected multiple times, because after you pick an item, you put it back before the next pick.
\sssc{Distinguishable objects choose with replacement}
The number of ways to choose with replacement $m$ objects out of $n$ distinguishable objects is
\[m^n.\]
\subsection{Binomial coefficients (二項次係數) or combinations (組合數)}
Below, we will define binomial coefficient $\binom{n}{m}$, also known as combination $C_m^n$, called $n$ choose $m$.
\subsubsection{Nonnegative integer}
The number of ways to choose $m$ objects out of $n$ distinguishable objects is
\[\frac{n!}{m!(n-m)!}.\]
For all $n,m\in\bbN_0$, define
\[\binom{n}{m}=\bcs
\frac{n!}{m!(n-m)!},\quad&m\le n\\
0,\quad&\tx{otherwise}
\ecs.\]
\subsubsection{Complex number $n$ and nonnegative integer $m$}
For all $n\in\bbC$ and $m\in\bbN_0$, define
\[\binom{n}{m}=\frac{\prod_{i=0}^{m-1}(n-i)}{m!}.\]
It satisfies
\[\binom{n}{m}=(-1)^m\cdot\binom{-n+m-1}{m}.\]
\subsubsection{$\Re(m),\Re(n-m)>-1$}
For all $n,m\in\bbC$ with $\Re(m),\Re(n-m)>-1$,
\[\binom{n}{m}=\frac{\Gamma(n+1)}{\Gamma(m+1)\Gamma(n-m+1)}.\]
\subsubsection{Complex number}
For all $n,m\in\bbC$,
\[\binom{n}{m}=\bcs
\frac{\Gamma(n+1)}{\Gamma(m+1)\Gamma(n-m+1)},\quad&\Re(m),\Re(n-m)>-1\\
(-1)^m\cdot\binom{-n+m-1}{m},\quad&\Re(n)\le -1\\
0,\quad&\tx{otherwise}
\ecs.\]
\subsubsection{Binomial theorem (二項式定理)}
For $n,x\in\bbC$, if either 
\bit
\item $|x|<1$, or
\item $\Re(n)>-1$ and $x=1$, or
\item $m\ge 0$ and $|x|=1$, or
\item $m\in\bbN_0$,
\eit
then:
\[\qty(1+x)^n=\sum_{m=0}^{\infty}\binom{n}{m}x^m,\]
where the RHS is called binomial series.
\bpr
PLACEHOLDER
\epr
\ssc{Multinomial coefficient}
\subsubsection{Nonnegative integer}
The number of ways to partition $n$ distinguishable objects into $m$ labeled groups with size $n_1,n_2,\ldots,n_m$ is
\[\frac{n!}{\prod_{i=1}^mn_i!}.\]
For all $n,m,n_1,n_2,\ldots,n_m\in\bbN_0$ with $\sum_{i=1}^mn_i=n$, define
\[\binom{n}{n_1,n_2,\ldots,n_m}=\bcs
\frac{n!}{\prod_{i=1}^mn_i!},\quad&m\le n\\
0,\quad&\tx{otherwise}
\ecs.\]
The number of ways to partition $k$ of $n$ distinguishable objects into $m$ labeled groups with size $n_1,n_2,\ldots,n_m$ is
\[\frac{(n)_k}{\prod_{i=1}^mn_i!}.\]
For all $n,m,n_1,n_2,\ldots,n_m\in\bbN_0$, let $k=\sum_{i=1}^mn_i\le n$, define
\[\binom{n}{n_1,n_2,\ldots,n_m}=\bcs
\frac{(n)_k}{\prod_{i=1}^mn_i!},\quad&m\le n\\
0,\quad&\tx{otherwise}
\ecs.\]
\subsubsection{Complex number $n$ and nonnegative integer $n_i$}
For all $n\in\bbC$ and $m,n_1,n_2,\ldots,n_m\in\bbN_0$, let $k=\sum_{i=1}^mn_i$, define
\[\binom{n}{n_1,n_2,\ldots,n_m}=\frac{(n)_k}{\prod_{i=1}^mn_i!}.\]
It satisfies
\[\binom{n}{n_1,n_2,\ldots,n_m}=(-1)^k\cdot\binom{-n+m-1}{n_1,n_2,\ldots,n_m}.\]
\subsubsection{$\Re(n_i),\Re(n-k)>-1$}
For all $m\in\bbN_0$ and $n,n_1,n_2,\ldots,n_m\in\bbC$ with $\Re(n_i),\Re(n-k)>-1$, let $k=\sum_{i=1}^mn_i$, define
\[\binom{n}{n_1,n_2,\ldots,n_m}=\frac{\Gamma(n+1)}{\prod_{i=1}^m\Gamma(n_i+1)\Gamma(n-k+1)}.\]
\subsubsection{Complex number}
For all $m\in\bbN_0$ and $n,n_1,n_2,\ldots,n_m\in\bbC$, let $k=\sum_{i=1}^mn_i$, define
\[\binom{n}{n_1,n_2,\ldots,n_m}=\bcs\frac{\Gamma(n+1)}{\prod_{i=1}^m\Gamma(n_i+1)\Gamma(n-k+1)},\quad&\Re(n_i),\Re(n-k)>-1\\
(-1)^k\cdot\binom{-n+m-1}{n_1,n_2,\ldots,n_m},\quad&\Re(n_i),\Re(-n+m-1-k)>-1\\
0,\quad&\tx{otherwise}
\ecs.\]
\subsubsection{Multinomial theorem (多項式定理)}
For $n,x_1,x_2,\ldots,x_m\in\bbC$, if $\sum_{i=1}^m|x_i|<1$, then:
\[\qty(\sum_{i=1}^mx_i)^n=\sum_{n_1,n_2,\ldots,n_m\ge 0}\binom{n}{n_1,n_2,\ldots,n_m}\prod_{i=1}^mx_i^{\pht{i}n_i},\]
where the RHS is called multinomial series.
\bpr
PLACEHOLDER
\epr
\ssc{組合恆等式}
\sssc{范德蒙恆等式（Vandermonde identity）}
\[\sum_{k=0}^n\binom{n}{k}^2=\binom{2n}{n}.\]
\begin{proof}\mbox{}\\
Consider a lattice path from \((0,0)\) to \((n,n)\) in a grid where each step moves either right \((1,0)\) or up \((0,1)\).\\
Since we must choose \( n \) steps to move right out of the total \( 2n \) steps, the total number of such paths is given by:
\[\binom{2n}{n}.\]
Now, let’s count the same paths differently. Choose an intermediate point \((k, n-k)\) at step \( n \). The number of ways to reach this point from \((0,0)\) using \( k \) right steps and \( n-k \) up steps is \( \binom{n}{k} \). From \((k, n-k)\) to \((n,n)\), we need \( n-k \) right steps and \( k \) up steps, which can be done in \( \binom{n}{k} \) ways. Summing over all possible \( k \), we obtain:
\[\sum_{k=0}^{n}\binom{n}{k}^2.\]
\end{proof}
\sssc{重複組合}
由$n$類相異物中，任取$m$個為一組之方法數，其中每類物品的個數均不小於$m$且可重複選取。記作$H^n_m$。
\[
\begin{aligned}
H_m^n &= C^{n+m-1}_m \\
&= C^{n+m-1}_{n-1}
\end{aligned}
\]
\sssc{巴斯卡公式（Pascal's formula）}
\[C^n_m=C^{n-1}_{m-1}+C^{n-1}_m\]
\sssc{組合恆等式}
\[
\begin{aligned}
\sum_{k=0}^n C_k^n &= 2^n \\
\sum_{k=0}^n (-1)^k \cdot C_k^n &= 0 \\
\sum_{k=0}^n 2^k \cdot C_k^n &= 3^n \\
\sum_{k=0}^n k \cdot C_k^n &= n \cdot 2^{n-1} \\
\sum_{k=0}^n k^2 \cdot C_k^n &= n \cdot (n-1) \cdot 2^{n-2} + n \cdot 2^{n-1} \\
\sum_{k=3}^n C_2^k &= n \cdot 2^{n-1} \\
\left(\prod_{i=0}^x (k-i)\right) \cdot C^n_k &= \left(\prod_{i=0}^x (n-i)\right) \cdot C^{n-k-1}_{k-x-1} \quad x < k\land x \in \mathbb{N} \\
\sum_{k=0}^{n} \left(\left(\prod_{i=0}^x (k-i)\right) \cdot C^n_k\right) &= \sum_{k=0}^{n} \left(\left(\prod_{i=0}^x (n-i)\right) \cdot C^{n-k-1}_{k-x-1}\right) \\
&= \frac{n!}{\left(n-x-1\right)!} \cdot 2^{n-x-1} \\
& \quad x < n\land x \in \mathbb{N}, \text{定義}\,\forall k<0: C_k^n=0
\end{aligned}
\]
\sssc{巴斯卡三角形（Pascal's triangle）/楊輝三角形}
令最上面一列為第$0$列，向下每列遞增$1$；每一列最左之數為第$0$個，向右每個遞增$1$。巴斯卡三角形之第$n$列第$m$個數定義為$C_m^n$。
\[
\begin{aligned}
\begin{array}{ccccccc}
& & & 1 & & & \\
& & 1 & & 1 & & \\
& 1 & & 2 & & 1 & \\
1 & & 3 & & 3 & & 1 \\
& \vdots & & & & \vdots &
\end{array}
\end{aligned}
\]
\subsection{卡特蘭數（Catalan number）}
\sssc{卡特蘭數（Catalan number）}
所有在$n \times n$格點中不越過對角線的單調路徑的個數，一個單調路徑從格點左下角出發，在格點右上角結束，每一步均為向上或向右，記作$C_n$。

\tb{遞迴式}：
\[
\begin{cases}
C_0 &= 1 \\
C_n &= \sum_{k=1}^n C_{k-1}\cdot C_{n-k}\quad\text{（當\protect\ }n>0\text{）}
\end{cases}
\]
\begin{proof}
令對角線為$y=x$，第一次到達對角線上時的位置為$(k, k)$，則第一次到達對角線前的單調路徑數為$C_k$，第一次到達對角線後的單調路徑數為$C_{n-k}$，故對$n>0$成立。$C_0=C_1=0$，亦成立。
\end{proof}

\tb{一般式}：
\[
\begin{aligned}
C_n &= C_n^{2n} - C_{n-1}^{2n} \\
&= \frac{1}{n+1} \cdot C_n^{2n}
\end{aligned}
\]
\begin{proof}
自$(0, 0)$至$(n, n)$的單調路徑數為$C^{2n}_n$。考慮這些路徑中不符合卡特蘭數定義者，其第一次跨越對角線$y=x$的點必在$y=x+1$上，將接下來的路徑對$y=x+1$鏡射，則終點$(n,n)$變為$(n-1, n+1)$。在此$(n-1) \times (n+1)$格點中自$(0, 0)$至$(n-1, n+1)$的單調路徑數為$C^{2n}_{n-1}$。故$C_n = C_n^{2n} - C_{n-1}^{2n} = \frac{1}{n+1} \cdot C_n^{2n}$。
\end{proof}
\subsubsection{非正方形格點的卡特蘭數}
所有在$n \times k$格點中不越過對角線的單調路徑的個數，一個單調路徑從格點左下角出發，在格點右上角結束，每一步均為向上或向右，記作$C_{n,k}$。

\tb{一般式}：
\[C_{n,k}=
\begin{cases}
C_k^{n+k}-C_{k-1}^{n+k} & \quad\text{（當\protect\ }n \geq k\text{）} \\
C_n^{n+k}-C_{n-1}^{n+k} & \quad\text{（當\protect\ }n \leq k\text{）}
\end{cases}
\]
\end{document}